% Generated by experimento/render_final_tables.py from audited R2/A1 data.
\begin{table}[htbp]
\centering
\caption{Recuperação observada após o comando de retomada, em segundos}
\label{tab:final-recuperacao}
\small
\setlength{\tabcolsep}{4pt}
\begin{tabular}{lrrrrr}
\hline
Condição & Var. & Primeiro commit & Estável: $n/5$ & Confirmação & Backlog completo \\ \hline
Banco parado & S & 4,92 $\pm$ 0,47 & 5/5 & 38,47 $\pm$ 0,57 & -- \\
Banco parado & A & 4,44 $\pm$ 1,08 & 1/5 & 38,71 & -- \\
Consumidor parado & A & 11,82 $\pm$ 3,59 & 5/5 & 41,63 $\pm$ 7,54 & 15,36 $\pm$ 5,00 \\
\hline
\end{tabular}
\fonte{Elaborado pelo autor a partir da campanha R2/A1 de 27/10/2026.}
\vspace{0.2cm}
\begin{minipage}{\textwidth}
\footnotesize
Primeiro commit: primeiro marco observado, não prova de primeiro commit real. Confirmação exige três janelas completas de 10 s sob entrada ativa. Média/desvio são condicionados aos casos observados: banco síncrono $n=5$, banco assíncrono $n=1$, consumidor $n=5$. Não se preenchem tempos nos quatro casos assíncronos sem sequência qualificante. Backlog síncrono não se aplica (zero aceites pendentes na retomada); no banco assíncrono, não houve persistência integral dos aceites anteriores por haver DLQ; consumidor teve backlog integralmente persistido em cinco casos.
\end{minipage}
\end{table}
